\chapter{Aritmética Modular}

\section{Introducción}
$a \equiv b \pmod{n}$ si $n$ divide a $a - b$.

Propiedades: reflexiva, simétrica, transitiva.

\section{Operaciones Modulares}
\[
a + b \equiv (a \bmod n) + (b \bmod n) \pmod{n},
\]
\[
a \cdot b \equiv (a \bmod n) \cdot (b \bmod n) \pmod{n},
\]
\[
a^b \bmod n.
\]

\section{Algoritmo de Euclides}

\begin{algorithm}[htbp]
\caption{Algoritmo de Euclides}
\begin{algorithmic}[1]
\Require $a \geq b > 0$
\Ensure $\gcd(a,b)$
\While{$b \neq 0$}
  \State $r \gets a \bmod b$
  \State $a \gets b$
  \State $b \gets r$
\EndWhile
\State \Return $a$
\end{algorithmic}
\end{algorithm}

El algoritmo extendido encuentra $x, y$ tales que $ax + by = \gcd(a,b)$.

\section{Inversos Modulares}
$a^{-1}$ existe si y solo si $\gcd(a,n) = 1$.

\section{Pequeño Teorema de Fermat}
Si $p$ es primo y $p \nmid a$:
\[
a^{p-1} \equiv 1 \pmod{p}.
\]

\section{Implementación en Haskell}

\begin{lstlisting}[caption={Aritmética modular en Haskell}]
mcd :: Integral a => a -> a -> a
mcd a 0 = abs a
mcd a b = mcd b (a `mod` b)

euclidesExtendido :: Integral a => a -> a -> (a, a, a)
euclidesExtendido a 0 = (a, 1, 0)
euclidesExtendido a b =
    let (g, x1, y1) = euclidesExtendido b (a `mod` b)
    in (g, y1, x1 - (a `div` b) * y1)

inversoModular :: Integral a => a -> a -> Maybe a
inversoModular a m =
    let (g, x, _) = euclidesExtendido a m
    in if g == 1 then Just (x `mod` m) else Nothing

potenciaMod :: Integer -> Integer -> Integer -> Integer
potenciaMod base 0 _ = 1
potenciaMod base exp mod
    | even exp = let p = potenciaMod base (exp `div` 2) mod
                 in (p*p) `mod` mod
    | otherwise = (base * potenciaMod base (exp - 1) mod) `mod` mod
\end{lstlisting}

\section{Ejercicios}
\begin{enumerate}
  \item Calcular el inverso modular de 3 módulo 11.
  \item Demostrar el Pequeño Teorema de Fermat.
  \item Implementar en Haskell una función que verifique si un número es primo usando el test de Fermat.
\end{enumerate}